% FilterDDP convergence for the two bounded cases.
%
% Data from figures/convergence_{6e_T6,6d_T3}.csv, written by
% examples/power_system/convergence_trace.jl. FilterDDP keeps no iteration
% history (src/solver_data.jl holds only the current iterate), so the trace is
% captured from its verbose stream and parsed back; see that script's header,
% including why the printed alpha/ls columns have to be de-staggered.
%
% B&W: each series carries colour AND a distinct dash pattern AND a distinct
% marker, so it survives a greyscale printout on shape alone.
%   pr_inf : blue  / solid     / circle
%   du_inf : green / dashed    / square
%   cs_inf : red   / dash-dot  / triangle
%   mu     : grey  / dotted    / none
%
\begin{figure}[!t]
\centering
\begin{tikzpicture}
\pgfplotsset{
  convaxis/.style={
    % NOTE: ymode/ymin/ymax are set per-axis below, not here. Carrying `ymode`
    % inside a \pgfplotsset style makes it resolve against /tikz/ at \end{axis}
    % and error out with "I do not know the key '/tikz/ymode'".
    width=0.78\columnwidth, height=0.42\columnwidth, scale only axis,
    tick align=outside,
    tick label style={font=\scriptsize},
    label style={font=\scriptsize},
    axis line style={black, line width=0.6pt},
    grid=both,
    major grid style={line width=0.6pt, color=gridMajor, opacity=0.7},
    minor grid style={line width=0.4pt, color=gridMinor, opacity=0.5},
    ylabel near ticks,
  },
  convseries/.style={line width=1.0pt, mark size=1.3pt},
}

% ---------------- Panel 1: 6e, T = 6 ----------------
\begin{axis}[
    convaxis, name=axA,
    ymode=log, ymin=1e-17, ymax=5e0, ytick={1e-16,1e-12,1e-8,1e-4,1e0},
    xmin=-0.5, xmax=16.5, xtick={0,4,8,12,16},
    xticklabels={},
    ylabel={$T = 6$ residual},
    legend style={at={(0.5,1.05)}, anchor=south, legend columns=4,
                  draw=none, fill=none, column sep=0.7ex, font=\scriptsize},
    legend cell align=left,
]
\addplot[convseries, cPr, mark=*, mark options={fill=cPr, draw=cPr}]
  table[col sep=comma, x=k, y=pr] {figures/convergence_6e_T6.csv};
\addlegendentry{$\mathrm{pr}_\infty$}
\addplot[convseries, cDu, densely dashed, mark=square*,
         mark options={fill=cDu, draw=cDu, solid}]
  table[col sep=comma, x=k, y=du] {figures/convergence_6e_T6.csv};
\addlegendentry{$\mathrm{du}_\infty$}
\addplot[convseries, cCs, dash dot, mark=triangle*,
         mark options={fill=cCs, draw=cCs, solid}]
  table[col sep=comma, x=k, y=cs] {figures/convergence_6e_T6.csv};
\addlegendentry{$\mathrm{cs}_\infty$}
\addplot[convseries, cMu, densely dotted, line width=1.2pt]
  table[col sep=comma, x=k, y=mu] {figures/convergence_6e_T6.csv};
\addlegendentry{$\mu$}
\end{axis}

% ---------------- Panel 2: 6d, T = 3 ----------------
\begin{axis}[
    convaxis, at={(axA.below south west)}, anchor=north west, yshift=-2.5mm,
    ymode=log, ymin=1e-17, ymax=5e0, ytick={1e-16,1e-12,1e-8,1e-4,1e0},
    xmin=-0.5, xmax=26.5, xtick={0,5,10,15,20,25},
    xlabel={Iteration $k$},
    ylabel={$T = 3$ residual},
]
\addplot[convseries, cPr, mark=*, mark options={fill=cPr, draw=cPr}]
  table[col sep=comma, x=k, y=pr] {figures/convergence_6d_T3.csv};
\addplot[convseries, cDu, densely dashed, mark=square*,
         mark options={fill=cDu, draw=cDu, solid}]
  table[col sep=comma, x=k, y=du] {figures/convergence_6d_T3.csv};
\addplot[convseries, cCs, dash dot, mark=triangle*,
         mark options={fill=cCs, draw=cCs, solid}]
  table[col sep=comma, x=k, y=cs] {figures/convergence_6d_T3.csv};
\addplot[convseries, cMu, densely dotted, line width=1.2pt]
  table[col sep=comma, x=k, y=mu] {figures/convergence_6d_T3.csv};
\end{axis}
\end{tikzpicture}
\caption{FilterDDP convergence, all bounds imposed: $T = 6$ (top, 17 iterations)
and $T = 3$ (bottom, 27).  Feasibility is reached almost immediately --- both
$\mathrm{pr}_\infty$ and $\mathrm{du}_\infty$ hit machine precision by $k = 4$
--- after which progress is entirely the barrier parameter $\mu$ driving
$\mathrm{cs}_\infty$ down alongside it.  Iteration count here is set by the
$\mu$ schedule, not by the optimisation.}
\label{fig:convergence}
\end{figure}
